\documentclass[11pt]{article}
\usepackage[margin=0.85in]{geometry}
\usepackage{amsmath,amssymb,hyperref}
\hypersetup{colorlinks=true,linkcolor=blue,urlcolor=blue}
\newcommand{\F}{\mathbb F}
\newcommand{\Z}{\mathbb Z}
\newcommand{\val}{v_p}
\title{\vspace{-2em}A sharp residue cutoff for spanning trees\\ in abelian Cayley towers}
\author{Research supplement to the spanning-tree method transfer}
\date{26 November 2026}
\begin{document}
\maketitle
\begin{abstract}
For odd \(p\), the steps modulo \(p\) determine the \(p\)-adic valuation of
the spanning-tree count at every level of a connected abelian Cayley tower
with fewer than \(p\) undirected step pairs. This improves the previous
sufficient bound \(2m<p\) to the sharp uniform bound \(m<p\).
At \(m=p\), two simple rank-two towers with identical reductions modulo \(p\)
give valuations of different growth orders: \(np^n\) and \(p^n\).
We give both valuations exactly at every level.
\end{abstract}

\section{Residue formula in arbitrary rank}
Let \(p\) be odd, let \(s_1,\ldots,s_m\in\Z^d\) span modulo \(p\), and put
an undirected edge from \(x\) to \(x+s_i\) for each \(x,i\) on
\((\Z/p^n\Z)^d\). Keep edge multiplicities; loops do not affect spanning
trees. Write \(\kappa_n\) for the count. For \(a\in\F_p^d\setminus\{0\}\),
define the following polynomial in \(\F_p[T]/(T^p)\):
\[
 L_a(T)=2m-\sum_{i=1}^m
 \big((1+T)^{[a\cdot s_i]}+(1+T)^{[-a\cdot s_i]}\big),
\]
where brackets denote representatives in \(\{0,\ldots,p-1\}\).
Suppose every \(L_a\) is nonzero, and let \(j(a)\) be its smallest nonzero
degree. Then
\begin{equation}
 \val(\kappa_n)=C\sum_{h=1}^n p^{(d-1)(h-1)}-dn,\qquad
 C=\frac{1}{p-1}\sum_{a\in\F_p^d\setminus\{0\}}j(a).
 \label{main}
\end{equation}
Thus only binomial coefficients modulo \(p\) are needed. The hypothesis is
equivalent to requiring that in every nonzero direction, some nonzero
unordered pair \(\{r,-r\}\) of projected residues occurs among the \(m\)
projected steps a number of times not divisible by \(p\).
Indeed, in the variable \(Z=1+T\) the coefficient at \(Z^r\) is the
negative of that multiplicity for \(r\ne0\); the coefficients sum to zero.
Spanning guarantees a nonzero projected residue in every nonzero
direction, so \(m<p\) suffices.

\paragraph{Cyclotomic evaluation.}
A character of exact order \(p^h\) has primitive vector
\(b\in(\Z/p^h\Z)^d\), meaning at least one coordinate is a unit.
Fix a primitive \(p^h\)-th root \(\zeta\), set \(\pi=\zeta-1\), and
\(e_h=(p-1)p^{h-1}\), so \(\val(\pi)=1/e_h\).
Its Laplacian eigenvalue is \(F_b(\zeta)\), where the integer polynomial
\(F_b(Z)\) is given by the formula for \(L_a\) but with \(Z\) in place of
\(1+T\) and representatives modulo \(p^h\).
In particular \(F_b(1)=0\).
\newpage
The coefficients of \(F_b(1+T)\) in degrees below \(p\), modulo \(p\), are
those of \(L_{\bar b}(T)\): over \(\F_p\), \((1+T)^p=1+T^p\).
Whenever \(j(\bar b)<e_h\), its term at \(\pi\) is uniquely of minimum
valuation \(j(\bar b)/e_h\). Smaller degrees have coefficients divisible
by \(p\); higher degrees have valuation at least their degree divided
by \(e_h\).

The only boundary case is \(h=1\), \(j=p-1\).
Then \(F_b(Z)\) has degree at most \(p-1\) and is congruent to
\(c\Phi_p(Z)\bmod p\) for an integer unit \(c\bmod p\), since
\(\Phi_p(Z)\equiv(Z-1)^{p-1}\bmod p\).
Write \(F_b=c\Phi_p+pG\). Evaluating at 1 gives \(G(1)=-c\).
As \(\zeta\equiv1\bmod\pi\), \(G(\zeta)\) is a unit and the valuation is
\(1=j/e_1\) as claimed. This uses the exact zero at 1 and does not require
an upper bound on the multiplicities. Each nonzero direction modulo \(p\)
has \(p^{d(h-1)}\) primitive lifts at order \(p^h\). The matrix-tree
theorem multiplies the eigenvalues of all nontrivial characters and
divides by the vertex count \(p^{dn}\). Summing valuations proves
\eqref{main}.

The same reduction works with more residue digits: if a primitive
projection \(b\) is known modulo \(p^k\), expand \(F_b(1+T)\) in
\(\F_p[T]/(T^{p^k})\) using representatives modulo \(p^k\).
If its order there is \(j<e_h\) and \(h\ge k\), all lifts to order
\(p^h\) give valuation \(j/e_h\). In particular, any nonzero such
polynomial suffices for \(h>k\). This is useful when \(L_{\bar b}=0\).

\section{Sharpness and two exact growth laws}
For \(0\le i<p\), take steps \(s_i=(1,i)\) and \(s_i'=(1+pi^2,i)\),
respectively. These give simple connected graphs at every level:
the steps and their negatives are distinct modulo \(p\), and
the first two steps span modulo \(p\). The graphs agree at level one.
Put \(E=p(p-1)\), \(\delta=\mathbf 1_{p\equiv3\bmod4}\). The exact laws
are
\begin{align}
 \val(\kappa_n)&=\sum_{j=0}^{n-1}
        \big[p^2-1+j(p-1)^2\big]p^j=((p-1)n+1)p^n-1, \label{flat}\\
 \val(\kappa_n')&=E+p+1+(3E/2+\delta)\sum_{j=1}^{n-1}p^j-2n. \label{curved}
\end{align}
Their difference at \(n=2\) is \(pE/2+2-p\delta>0\). So even for simple
rank-two graphs the strict bound \(p>m\) cannot be relaxed to \(p\ge m\).
The growth orders differ, not just the initial values.

\paragraph{Common contributions.}
At level one, an order-\(p\) character with vector \((a,b)\) has
eigenvalue \(2p\) if \(b\ne0\), since the projections run through all
residues; otherwise the eigenvalue is \(p(2-\zeta^a-\zeta^{-a})\)
with valuation \(1+2/(p-1)\). The layer contribution before subtracting
the vertex-count term is \(E+p+1\).
For \(h\ge2\) and \(b\) a unit, the mod-\(p\) polynomial is
\(-2\sum_{r=0}^{p-1}(1+T)^r=-2T^{p-1}\).
Each such eigenvalue has valuation \((p-1)/e_h\), contributing
\(E p^{h-1}\) in total. The remaining primitive vectors have \(a\) a
unit and \(p\mid b\).
\newpage
\paragraph{Straight steps.}
If \(b=0\) at order \(p^h\), the contribution is \(e_h+2\).
When \(t=\val(b)\) satisfies \(1\le t\le h-1\), the eigenvalue is
\[
 2p-\big(\zeta^a+\zeta^{-a-(p-1)b}\big)
       \sum_{i=0}^{p-1}\zeta^{bi}.
\]
The parenthesis is a unit because it reduces to \(2\bmod\pi\).
If \(t<h-1\), the geometric sum has valuation
\((p-1)p^t/e_h<1\), by taking the ratio
\((1-\zeta^{pb})/(1-\zeta^b)\).
For \(t=h-1\) the geometric sum is zero and the eigenvalue has
valuation 1, still \((p-1)p^t/e_h\). There are \(e_h\) choices for \(a\)
and \((p-1)p^{h-1-t}\) choices for \(b\) at each \(t\).
Each \(t\) therefore contributes \((p-1)^2p^{h-1}\).
Adding the shared cases and subtracting \(2n\) gives \eqref{flat}.

\paragraph{Quadratic steps.}
Write \(a=a_0+pa_1\) with \(a_0\in\{1,\ldots,p-1\}\), and \(b=pc\).
The positive projected residues are congruent modulo \(p^2\) to
\(a_0+p k_i\), where \(k_i=a_1+ci+a_0i^2\bmod p\).
With \(U=T^p\) and \(Z=1+T\), the sum of positive projected terms
modulo \((p,T^{p^2})\) is \(Z^{a_0}\sum_i(1+U)^{k_i}\).
The negative sum has \(Z^{p-a_0}\) and exponents \(-1-k_i\bmod p\).
Let \(r=(p-1)/2\). For each \(j<r\) the coefficient of \(U^j\)
vanishes in both sums: the polynomial in \(i\) has degree \(2j<p-1\),
and every monomial in that range sums to zero over \(\F_p\), including
the constant. At \(j=r\) the positive coefficient is
\(A=-a_0^r/r!\ne0\), from the degree \(p-1\) term, and the
negative coefficient is \((-1)^r A\).

Consequently, the first surviving degree in the Laplacian is \(pr\)
if \(r\) is even. If \(r\) is odd, the constant of the \(Z\)-factor
cancels, but its \(T\)-coefficient is \(2a_0 A\) up to the
overall Laplacian sign and is nonzero. The degree is then \(pr+1\).
In either case the degree is \(E/2+\delta<e_h\). Dividing by \(e_h\)
gives the eigenvalue valuation. There are \(e_h p^{h-1}\) such vectors,
so the common and quadratic contributions sum to
\((3E/2+\delta)p^{h-1}\) for \(h\ge2\).
The matrix-tree formula gives \eqref{curved}.

\section*{Sources and verification}
This supplements the \href{https://chatgpt.com/c/6b07d086-1ef0-83e8-9354-9bbcf6afcaf3?messageId=88092bdb-6383-595f-a45a-d402786a8c09}{earlier fourteen-formula report}.
The valuation step is adapted from the
\href{https://github.com/lsernau/grh-lean/blob/72fead102f663d5c7a407964a13dc0cf9d977207/grh/src/Lib/PadicHodge/TateSen.lean\#L131-L157}{pinned Tate--Sen argument};
neither RH nor GRH is an input. The comparison results are the
general growth theorem of \href{https://alco.centre-mersenne.org/articles/10.5802/alco.304/}{DuBose--Valli\`eres}
and \href{https://arxiv.org/abs/2606.03579}{Kataoka's asymptotic analysis}.
The companion script checks the formulas by exact cyclotomic polynomial
reduction and, for small graphs, by independent integer cofactors.
At \(p=3,n=2\), those \(80\times80\) cofactors give \(44\) and \(36\).
This is a written proof, not a Lean formalization or a claim of priority.
\end{document}
